\documentclass[12pt]{article}
\usepackage[utf8]{inputenc}
\usepackage[margin=1in]{geometry}
\usepackage{amsmath}
\usepackage{graphicx}
\usepackage{booktabs}
\usepackage{hyperref}
\usepackage[round]{natbib}

\title{Hippocampal and Striatal Volume Contributions to the Longitudinal Development of Feedback-Timing-Dependent Reinforcement Learning in 6-to-7-Year-Old Children}
\author{Author A$^{1}$, Author B$^{1,2}$, Author C$^{2}$, Author D$^{1}$ \\
\small $^{1}$Department of Developmental Psychology, University of Amsterdam, Amsterdam, Netherlands \\
\small $^{2}$Amsterdam Brain and Cognition Center, University of Amsterdam, Amsterdam, Netherlands}
\date{}

\begin{document}
\maketitle

\begin{abstract}
Adult studies demonstrate that feedback timing modulates engagement of hippocampal (delayed feedback) versus striatal (immediate feedback) memory systems during reinforcement learning, but whether similar dissociations exist in developing children is unknown. Here, we report a two-year longitudinal study of 142 children (wave 1 age: mean 7.19, SD 0.46 years) who completed a probabilistic reinforcement learning task with immediate and delayed feedback conditions. Children showed significant longitudinal improvements: increased accuracy ($\beta = 0.550$, $p < 0.001$), increased win-stay behavior ($\beta = 0.586$, $p < 0.001$), decreased lose-shift behavior ($\beta = -0.252$, $p < 0.001$), and faster reaction times ($\beta = -221$ ms, $p < 0.001$). Computational modeling identified the best-fitting model as one with a shared learning rate and feedback-timing-dependent inverse temperature. Learning rates (0.02--0.05) were far below simulated adult optima ($\sim$0.29), with both learning rate and inverse temperature increasing toward more optimal parameter combinations. In a subsample with longitudinal MRI ($n = 73$), hippocampal volume showed more protracted maturation than striatal volume. Latent change score models revealed that larger hippocampal volume was associated with better delayed learning, consistent with adult findings. Unexpectedly, larger striatal volume was associated with both immediate and delayed learning, suggesting a less differentiated, more cooperative role for memory systems in middle childhood compared to adults.
\end{abstract}

\section{Introduction}

Reinforcement learning (RL)---the process of learning from outcomes to guide future decisions---is a fundamental cognitive capacity that develops throughout childhood and adolescence \citep{sutton1998rl, nussenbaum2020development}. In adults, value-based RL engages dissociable memory systems: the striatum supports learning from immediate feedback through incremental stimulus-response associations, while the hippocampus supports learning when feedback is delayed, leveraging episodic memory to bridge the temporal gap between choice and outcome \citep{foerde2006hippocampus, shohamy2008striatum}.

This dual-system framework predicts that the developmental trajectories of hippocampal and striatal function should differentially influence how children learn from immediate versus delayed feedback. However, developmental studies of RL have been predominantly cross-sectional and have not incorporated feedback timing manipulations or computational models that decompose learning into component processes \citep{palminteri2016computational}. Furthermore, the hippocampus and striatum follow different maturational timelines, with the hippocampus showing more protracted structural development \citep{goddings2014puberty}, yet how these structural differences relate to RL parameter changes over time in children has not been tested.

Here, we address these gaps with a two-year longitudinal study combining behavioral RL assessment, computational modeling with Bayesian model comparison, structural MRI, and latent change score modeling. We test three key questions: (1) How does value-based learning improve longitudinally in middle childhood? (2) Does feedback timing modulate learning parameters and their development? (3) Do hippocampal and striatal volumes differentially predict learning changes as a function of feedback timing?

\section{Methods}

\subsection{Participants}

One hundred forty-two children were enrolled at wave 1 (46\% female; mean age 7.19 $\pm$ 0.46 years), with 127 returning at wave 2 (mean age 9.25 $\pm$ 0.45 years; mean inter-wave interval 2.07 $\pm$ 0.17 years; Table~\ref{tab:participants}). A subset underwent structural MRI at both waves ($n = 73$ with longitudinal MRI data).

\begin{table}[htbp]
\centering
\caption{Participant characteristics.}
\label{tab:participants}
\begin{tabular}{lcc}
\toprule
Parameter & Wave 1 & Wave 2 \\
\midrule
Enrolled & 142 & 127 \\
Included in RL analysis & 140 & 126 \\
Sex (\% female) & 46\% & 46\% \\
Age mean (years) & 7.19 & 9.25 \\
Age SD (years) & 0.46 & 0.45 \\
MRI scanned & 90 & 104 \\
Usable structural scans & 82 & 99 \\
Longitudinal MRI & \multicolumn{2}{c}{73} \\
\bottomrule
\end{tabular}
\end{table}

\subsection{Reinforcement Learning Task}

Children completed a probabilistic RL task with four cue-choice pairs (87.5\% contingent reward probability) across two feedback timing conditions. In the immediate condition, feedback followed choice after 1 second. In the delayed condition, a 6-second delay with an incidentally presented object preceded feedback.

\subsection{Behavioral Analysis}

Accuracy, win-stay probability, lose-shift probability, and reaction time were analyzed with generalized linear mixed models (GLMMs) with wave and feedback condition as predictors.

\subsection{Computational Modeling}

Multiple value-based RL models were compared using Bayesian model comparison (Pseudo-BMA+ with Bayesian bootstrap, 100,000 iterations) \citep{yao2017pseudobma, daw2011trial}. The winning model (vbm3) estimated a shared learning rate ($\alpha$) and feedback-timing-dependent inverse temperature ($\tau_\text{immediate}$, $\tau_\text{delayed}$). Parameter and model recovery analyses confirmed reliability.

\subsection{MRI and Brain-Cognition Modeling}

Hippocampal and striatal volumes were extracted from T1-weighted structural MRI. Brain-cognition associations were modeled using four-variate latent change score (LCS) models capturing longitudinal change in striatal volume, hippocampal volume, immediate learning score, and delayed learning score \citep{mcewen2021latent}.

\section{Results}

\subsection{Longitudinal Behavioral Improvements}

Children showed robust longitudinal improvements across all behavioral measures (Table~\ref{tab:glmm}). Accuracy increased significantly between waves, win-stay behavior increased (indicating better exploitation of positive feedback), and lose-shift behavior decreased (indicating reduced sensitivity to negative outcomes). Reaction times were faster at wave 2 and faster for delayed versus immediate feedback.

\begin{table}[htbp]
\centering
\caption{GLMM fixed effects for behavioral variables.}
\label{tab:glmm}
\begin{tabular}{llcccc}
\toprule
DV & Predictor & $\beta$ & SE & Statistic & $p$ \\
\midrule
Accuracy & Wave 2 & 0.550 & 0.061 & $z = 8.97$ & $< 0.001$ \\
Accuracy & Delayed & 0.013 & 0.024 & $z = 0.54$ & 0.590 \\
Win-Stay & Wave 2 & 0.586 & 0.071 & $z = 8.22$ & $< 0.001$ \\
Win-Stay & Delayed & 0.023 & 0.033 & $z = 0.69$ & 0.489 \\
Lose-Shift & Wave 2 & $-0.252$ & 0.037 & $z = -6.87$ & $< 0.001$ \\
RT (ms) & Wave 2 & $-221$ & 22.8 & $t = -9.70$ & $< 0.001$ \\
RT (ms) & Delayed & $-13.8$ & 6.59 & $t = -2.10$ & 0.038 \\
\bottomrule
\end{tabular}
\end{table}

Feedback timing did not significantly modulate accuracy, win-stay, or lose-shift behavior. However, delayed feedback was associated with significantly faster reaction times ($\beta = -13.8$ ms, $p = 0.038$).

\subsection{Computational Modeling: Feedback-Timing-Dependent Inverse Temperature}

Bayesian model comparison selected vbm3 (shared learning rate, separate inverse temperature by feedback timing) as the best-fitting model at both group and individual levels. Both the learning rate and inverse temperature increased between waves, reflecting movement toward more optimal parameter combinations (Table~\ref{tab:params}). Children's learning rates (0.02--0.05) remained far below simulated adult optima ($\sim$0.29), indicating substantial room for continued development.

\begin{table}[htbp]
\centering
\caption{Winning model parameter estimates across waves.}
\label{tab:params}
\begin{tabular}{lccc}
\toprule
Parameter & Wave 1 (approx.) & Wave 2 (approx.) & Direction \\
\midrule
$\alpha$ (learning rate) & 0.02 & 0.05 & Increased \\
$\tau_\text{immediate}$ & Lower & Higher & Increased \\
$\tau_\text{delayed}$ & Lower & Higher & Increased \\
Simulated optimal $\alpha$ & \multicolumn{2}{c}{$\sim$0.29} & -- \\
Optimal accuracy & \multicolumn{2}{c}{96.5\%} & -- \\
\bottomrule
\end{tabular}
\end{table}

\subsection{Brain Volume Development}

Both hippocampal and striatal volumes increased between waves, but hippocampal volume showed more protracted maturation (larger relative increase), consistent with the known developmental sequence of subcortical structures \citep{goddings2014puberty}.

\subsection{Brain-Cognition Associations: Latent Change Score Models}

Four-variate LCS models revealed differential brain-cognition associations by feedback timing (Table~\ref{tab:lcs}). Larger hippocampal volume was associated with better delayed learning longitudinally, consistent with the adult literature supporting hippocampal involvement in delayed-feedback learning \citep{foerde2006hippocampus}. However, larger striatal volume was associated with both immediate and delayed learning, contrary to the adult prediction of striatal specificity for immediate feedback \citep{frank2004carrot, shohamy2008striatum}.

\begin{table}[htbp]
\centering
\caption{Brain-cognition associations from latent change score models.}
\label{tab:lcs}
\begin{tabular}{llll}
\toprule
Volume & Learning Type & Association & Adult Prediction \\
\midrule
Hippocampal & Delayed & Positive (larger $=$ better) & Consistent \\
Hippocampal & Immediate & Weaker/NS & Consistent \\
Striatal & Immediate & Positive & Consistent \\
Striatal & Delayed & Positive & Inconsistent \\
\bottomrule
\end{tabular}
\end{table}

\subsection{Sensitivity Analyses}

Analyses excluding poor learners ($<$50\% accuracy in last 20 trials; wave 1: $n = 13$; wave 2: $n = 6$) yielded unchanged results, confirming that the findings are not driven by non-learning participants.

\section{Discussion}

This study provides the first longitudinal evidence for the development of value-based reinforcement learning in middle childhood, combining behavioral analysis, computational modeling, structural neuroimaging, and latent change score modeling. Three main findings emerge.

\subsection{Developmental Trajectory of Reinforcement Learning}

Children showed robust improvements across all behavioral and computational measures over the two-year interval. The increase in both learning rate and inverse temperature between waves indicates that children are becoming both better learners (updating values more effectively) and better decision-makers (choosing more consistently based on learned values) \citep{palminteri2016computational, nussenbaum2020development}. However, children's learning rates remained far below adult optima, indicating that the development of RL is far from complete by age 9--10.

The finding that the best-fitting model included feedback-timing-dependent inverse temperature but not learning rate suggests that the way children exploit learned values, rather than how quickly they learn, is differentially sensitive to feedback timing during this developmental period \citep{daw2011trial}.

\subsection{Less Differentiated Memory System Roles in Childhood}

The most striking finding is the non-specific association between striatal volume and both immediate and delayed learning. In adults, the striatum is preferentially engaged during immediate-feedback learning, with the hippocampus taking over for delayed-feedback conditions \citep{foerde2006hippocampus, shohamy2008striatum}. The broader striatal involvement observed here suggests that in middle childhood, the memory systems operate in a more cooperative, less differentiated manner \citep{mccormick2019hippocampal}. The hippocampal association with delayed learning is already present, but the striatal specialization for immediate feedback has not yet fully emerged. This developmental pattern is consistent with the protracted maturation of hippocampal structure relative to the striatum and may reflect the gradual differentiation of memory systems during middle childhood.

\subsection{Limitations}

The sample size for longitudinal MRI ($n = 73$) limits statistical power for detecting brain-cognition associations. The 87.5\% contingent reward probability may have limited variability in learning performance. The two-year longitudinal design captures a specific developmental window; longer follow-up would be needed to track the emergence of adult-like memory system specialization.

\section{Conclusion}

Middle childhood is characterized by substantial improvements in value-based reinforcement learning, driven by increases in both learning rate and value-guided decision-making. Feedback timing modulates the inverse temperature parameter but not accuracy. Brain-cognition associations reveal hippocampal involvement in delayed learning consistent with adult patterns, but a more cooperative, less differentiated role for the striatum than predicted by adult dual-system models. These findings suggest that the specialization of hippocampal and striatal memory systems for feedback-timing-dependent learning continues to develop beyond middle childhood.

\bibliographystyle{plainnat}
\bibliography{references}

\end{document}
